
\section{RQ2: In what ways are LLMs employed to secure SSC?}
\label{sec:rq2}

\par This Research Question examines how LLM defend Software Supply Chain
security, by means of reviewing collected literature. We begin by 
introducing the major attack vectors of SSC breaches, and then review
how LLMs are leveraged to mitigate these kinds of attacks. 

% IMPORTANT NOTE: this is not my original work.
\paragraph*{What are the attack vectors against SSC?} L. Williams et al.
categorized the three major attack vectors as follows\citep{sscdirection}:


% these are my words, i.e. paraphrased.
\begin{itemize}
  \item (Vector 1) Through malwares or vulnerabilities injected into open-source packages/libraries
  and other third-party dependencies (e.g. the \texttt{xz-util} incident 
  \citep{xzincident});

  \item (Vector 2) Through compromising the build infrastructure (e.g. CI/CD) where the
  software component is tested, packed and deployed;

  \item (Vector 3) Through targeting human developers involved in software development,
  such as \textit{Social Engineering}.
\end{itemize}

% Generated from dataset/rq2.json
\begin{figure}[ht]
  \centering
  \subcaptionbox{Distribution by year.}
    {\includegraphics[width=0.45\textwidth]{logo/rq2trend.pdf}}
  \subcaptionbox{Distribution by attack vectors.}
    {\includegraphics[width=0.45\textwidth]{logo/rq2dist.pdf}}

  \caption{Distribution of LLM-based Mitigation Publications.}
  \label{fig:rq2dist}
\end{figure}

\paragraph*{Trend Analysis.} The distributions of publications over
different types of attack vectors and year are shown in Figure~\ref{fig:rq2dist}.
We observe that defending Software Supply Chain against \textit{attack vector 1} is
steadily on the rise, with only few exceptions where defense to \textit{attack vector 2}
is studied. Apart from a few surveys in which \textit{attack vector 3} is introduced
(e.g. \cite{sscdirection}), we have found no studies proposing defense to
this kind of SSC attack.

% Generated from dataset/domain.json
% "package manager" : e.g. npm, maven, PyPI, etc.
% "AI Model" (platform) e.g. HuggingFace and application (e.g. LLM-integrated app
% \cite{10.1145/3756681.3757083})
% "Agent harness" includes Skills, plugins, MCP, etc.
\paragraph*{Which components of Software Supply Chain are secured by LLM?}

\subsection{Mitigation of Malware and Vulnerabilities}\label{subsec:vec1}


\begin{figure}[ht]
  \centering
  \includegraphics[width=\textwidth]{logo/vec1taxo.pdf}

  \caption{Taxonomy of Attack Vector 1 Mitigation Techniques. \\
   {\small \textit{det.} is short for the word detection; \textit{lib.} is short for the word library.} \\
   {\small Enclosure of circles implies logical enclosure; the size of circle
   is roughly proportional to \# studies obtained.}
  }
\end{figure}

\subsection{Mitigation of Build System Vulnerabilities}\label{subsec:vec2}


